% =====================================================================================
%  Lecture 4  —  23 September 2026   (guest lecture by Anson Ho, Epoch AI)
% =====================================================================================
\chapter{Will AI drive a software intelligence explosion?}\label{ch:L4}
\setlecturevideo{7AaNkZuYfmk}
\lectureinfoguest{23 September 2026}{55 min}{7AaNkZuYfmk}{Anson Ho (Epoch AI)}{recursive self-improvement, a growth-theory model, compute bottlenecks}

\begin{glance}
\begin{itemize}
  \item \textbf{The question.} Frontier labs say they want to automate AI research itself. If they succeed, does progress start to feed on itself, and can it do so with \emph{software alone}, with no new chips or data centers (a \term{software intelligence explosion})?
  \item \textbf{The basic model.} The Jones model of economic growth, $\dot S/S=R^{\lambda}S^{-\beta}$. When AIs do the research, $R\propto S$ and the loop blows up in finite time exactly when $\lambda>\beta$, i.e.\ (for $\beta>0$) when the \emph{returns to software R\&D} $r=\lambda/\beta$ exceed~$1$.
  \item \textbf{Estimating $r$.} Measure software quality as \emph{compute efficiency}; in a balanced-growth approximation $r=g_S/g_R$, the ratio of two growth rates. The data are poor: the estimates cannot rule $r>1$ in or out.
  \item \textbf{Extensions.} Experiments need compute, modeled by a CES production function whose substitution parameter decides whether compute is a \emph{bottleneck}; the number of AIs versus their intelligence; progress that depends on the scale of training; data from inside the labs.
  \item \textbf{Discussion.} The speaker's credence (about 20--25\% for an explosion), chips as a physical bottleneck, and why measuring capabilities is itself an unsolved problem.
\end{itemize}
\end{glance}

\begin{caution}[About the source for this lecture]
The recording has English auto-captions only, and the slides are not visible in them, so every figure I describe is described from the speaker's words. The captions garble several names (the speaker's own appears as ``Anen Paul'' in the first minute); those I could identify with confidence are written normally, the others are marked\flag. Claims about recent events are reported \emph{as made in the lecture}; where I could check one against a source, the check is in a ``Beyond the lecture'' box.
\end{caution}

\begin{figure}[tb]
\centering
\mapbar{3335}{600}{%
  \mapseg{0}{438}{slate}\mapseg{438}{807}{teal}\mapseg{807}{1462}{accent}\mapseg{1462}{2085}{amber}
  \mapseg{2085}{2558}{moss}\mapseg{2558}{3335}{plum}
  \mapin{219}{4.1}\mapin{622}{4.2}\mapin{1134}{4.3}\mapin{1773}{4.4.1 \ CES}\mapin{2321}{4.4.2--4}\mapin{2946}{4.5 \ Questions}
  \maplabel{219}{Race to automate}{3.2}{south west}
  \maplabel{622}{Jones model}{7.2}{south west}
  \maplabel{1134}{Measuring, estimating $r$}{3.2}{south west}
  \maplabel{1773}{Compute bottlenecks}{7.2}{south west}
  \maplabel{2321}{Beyond the basic model}{3.2}{south west}
}
\caption{\textbf{Lecture map} (55\,min, to scale). The talk builds the basic model (4.2), asks how it could be fitted to data (4.3), and spends its second half on what the basic model leaves out (4.4), before taking questions.}
\end{figure}

%% ------------------------------------------------------------------------------------
\section{The race to automate AI research}\label{sec:L4-race}
\tsm{0:01}
After a short introduction by the instructor, the guest, Anson Ho\flag{} of Epoch AI, says where his interest comes from. His work is to track what the frontier labs (mostly in Silicon Valley) are doing and how they try to build frontier AI. The talk is about one scenario: \emph{what happens if they succeed at their stated vision of automating AI research}, and thereby cause something like a \term{software intelligence explosion}?

\paragraph{The stated goal.}\tsm{1:04}
The starting point is what the people at the labs say they are doing. The speaker shows a post by Sam Altman, the CEO of OpenAI: the company's goal is to automate AI research \emph{fully} by March 2028. An earlier goal was a research \emph{intern} by September 2026, which the company now claims to have reached in a recent blog post, while it is ``racing toward'' automating the remaining tasks. (There has already been much progress in automating mathematics, as the class knows, he adds.)
The intended mechanism is \term{recursive self-improvement}: AIs develop better AIs, which develop better AIs still, because they improve their own algorithms and data. In the most extreme case this happens \emph{only in software}: no new chips, no new data centers, only the existing stock of physical resources, while the intelligence and the \emph{number} of AI systems grow dramatically. ``At least that is what they claim might happen.''

\begin{beyond}[What OpenAI actually published]
The goals quoted by the speaker are those Altman announced in October 2025: an ``intern-level'' AI research assistant by September 2026 and a ``legitimate AI researcher'' by March 2028. On 6~September 2026 OpenAI published a post, \emph{Research acceleration: the view inside OpenAI}\,\citep{openai2026ra}, declaring the first milestone reached \emph{by its own measurements}: an intern is defined there as a system that can carry out well-defined research tasks under human direction, including tasks that would take a skilled researcher a few days, and the post reports that, as of mid-August, its research organization ran about 3.1 agent-workdays of effort per workday of human labor. The page refused automated access, so these details are from search excerpts and press coverage, not from the post itself; the milestone is the company's own assessment and, as far as I could find, has not been verified independently.
\end{beyond}

\paragraph{How fast could it be?}\tsm{2:40}
The speaker's example of such a scenario is the \emph{AI Futures} model\,\citep{aifutures2025,aifutures2025dec}, in which the feedback loop compresses many years of progress into a few: over the last decade or so we went from simple vision models that tell cats from dogs to systems that, he says, contribute to settling the Navier--Stokes problem, and in the model one could see something like fifteen years of that kind of progress compressed into about three\flag. The stakes, he says, are very high, in three directions:
\begin{itemize}
  \item \textbf{Upside.} An enormous potential for scientific and medical innovation.
  \item \textbf{Misalignment.} The systems might not do what we care about; ``which is probably the motivation for this course''.
  \item \textbf{Concentration of power.} Imagine millions of AIs in the hands of a small number of actors, aligned to \emph{their} wishes. Even if the alignment problem were solved, our institutions might not be able to sustain such a concentration.
\end{itemize}

\begin{beyond}[Two results mentioned in passing]
The speaker twice mentions recent mathematical results by AI systems (``contribute to solving the Navier--Stokes problem'', and, later, the Jacobian conjecture). As far as I could check: on 8~September 2026 OpenAI announced that an internal system had produced a proof that the three-dimensional Navier--Stokes equations, driven by a smooth compactly supported \emph{force} and started from rest, can develop a singularity in finite time, with a Lean formalization\,\citep{openai2026ns}. That concerns the \emph{forced} alternative in the Clay problem statement; the unforced problem is not settled by it, and I found no independent refereed verification yet, so ``solved the Navier--Stokes problem'' says more than is established. For the Jacobian conjecture, press reports from July 2026 describe a counterexample found with AI assistance and checked by computer algebra; I did not check it against a primary source. (Lecture~1 discussed an earlier case, \citealp{openai2026unit}.)
\end{beyond}

\begin{discussion}[What is on the vertical axis?]
A student asks about the forecast plot, in which the line bends upward: is that line straight or exponential? The speaker: the vertical axis is the \emph{Epoch Capabilities Index} \citep{epoch2025eci}, whose meaning is not obvious. The points are real data; the extrapolation into the future is the AI Futures model's projection. One way to read the index is through the METR time-horizon graph (Lecture~1), with which it correlates very strongly, so that a level of the index corresponds to a task length (say, a one-day horizon). The observed trend in time horizon is an \emph{exponential}; the projection would be \emph{super-exponential}.
\end{discussion}

\paragraph{Pace or pause?}\tsm{5:43}
Given such stakes, the speaker says, employees and leaders of frontier companies have recently put out statements about \emph{pacing} the frontier, potentially even pausing or slowing down AI development. He names the chief executives of Anthropic, OpenAI and DeepMind and Elon Musk, and calls the position of Mark Zuckerberg ``a bit more unclear''; of the latter he thinks that he says that, in practice, many will naturally want to slow down while they solve alignment problems, because the cost of alignment is \emph{internalized} by whoever builds the systems\flag. The debate also reaches beyond tech: politicians (he names Bernie Sanders and Donald Trump) are weighing in, with different views. I have not verified these attributions.

\paragraph{What the talk will do.}\tsm{6:15}
The big question: \emph{if we reach the point of automating AI research, will advanced AI actually drive a software intelligence explosion?} The notion starts out vague (what does it mean for AI to improve its own software? how do we measure the capabilities of these systems?), so the aim is to make it rigorous: to explain the earliest attempts at a model, what the empirical evidence says about that model, and some extensions. ``Unfortunately, our current understanding is terrible; hopefully some of you can help improve on this state of the art, which is not good.''

%% ------------------------------------------------------------------------------------
\section{The basic model: Jones' growth theory}\label{sec:L4-jones}
\tsm{7:18}
To capture the feedback loop the speaker takes ``the most standard model of innovation out there'': the Jones model from the theory of economic growth\,\citep{jones1995}. (By a show of hands, some of the class know it; ``it is not actually that complicated'', and, he says, the mathematicians will know it better than he does.)

\subsection{Software quality and research effort}\label{sec:L4-setup}\tsm{7:48}
Let $S(t)>0$ be the \emph{quality of software}, one number for the moment (what it means is the topic of Section~\ref{sec:L4-measure}), and $R(t)\ge0$ the \emph{research input}: a production function that turns research effort into software quality. The model says that the \emph{growth rate} of $S$ depends on both,
\begin{equation}\label{eq:L4-jones}
   \frac{\dot S}{S}\;=\;R^{\lambda}\,S^{-\beta}
   \qquad\Big(\text{equivalently }\ \dot S=R^{\lambda}S^{\,1-\beta}\Big),
\end{equation}
up to a multiplicative constant that can be absorbed into the units of $R$. The two exponents have names:
\begin{description}
  \item[$\lambda$: the parallelization penalty.] If $\lambda<1$, ten times more researchers does not give ten times the growth rate. Ten times more researchers may give only about three times the benefit (for instance $\lambda=\tfrac12$ gives $\sqrt{10}\approx3$; the example with $\lambda$ is mine): ``people step on each other's toes''.
  \item[$\beta$: ideas get harder to find.] For $\beta>0$ the growth rate falls as the stock $S$ grows: there is less low-hanging fruit to pick. Depending on the sign of $\beta$ one could instead get a \emph{standing-on-the-shoulders-of-giants} effect (Newton), in which more existing knowledge makes new progress \emph{easier}.
\end{description}
A student asks whether the exponent lies between zero and one; ``I think zero to one''. (The captions do not make clear which exponent is meant\flag; the usual range of the parallelization exponent is $0<\lambda\le1$, and $\beta$ may have either sign.)

\subsection{Closing the loop}\label{sec:L4-loop}\tsm{9:50}
Normally the research input comes from human labor. In an intelligence explosion \emph{all} of it comes from AIs, so one replaces $R$ by something that depends on the quality of the software and on the amount of computation available to run the AI systems: ``I just mean we are plugging AI laborers into this production function''. In the simplest version,
\begin{equation}\label{eq:L4-RSC}
   R\;=\;S\cdot C ,
\end{equation}
where $C$ is the computation used to run the AI researchers (constant for now). Closing the loop and ``messing around with the symbols'' gives
\begin{equation}\label{eq:L4-closed}
   \dot S\;=\;S\,(SC)^{\lambda}S^{-\beta}\;=\;C^{\lambda}\,S^{\,1+\lambda-\beta}.
\end{equation}
Growth is \emph{hyperbolic}, meaning that the solution becomes infinite in finite time, exactly when the exponent of $S$ exceeds one, i.e.\ when $\lambda-\beta>0$, or, as the speaker prefers to write it,
\begin{equation}\label{eq:L4-r}
   r\;:=\;\frac{\lambda}{\beta}\;>\;1 .
\end{equation}
He calls $r$ the \term{returns to AI software R\&D}. The same parameter has been estimated for the economy as a whole in the well-known paper ``Are ideas getting harder to find?''\,\citep{bloom2020}: for total factor productivity its value is usually about $0.3$. Humans keep innovating, but new ideas get harder to find. The value of $r$ is the outcome of a \emph{race} between finding new ideas and running into diminishing returns, and it decides whether one gets hyperbolic growth (a software intelligence explosion) or a fizzle. ``We just simplified the whole thing into a single number, which is nice.''

\begin{figure}[tb]
\centering
\begin{tikzpicture}[>={Latex[length=2.2mm,width=1.6mm]}]
  \node[card=accent, minimum width=36mm, minimum height=13mm] (S) at (0,0) {software quality\\[-1pt]$S$};
  \node[card=teal, minimum width=36mm, minimum height=13mm] (R) at (6.8,0) {AI research labor\\[-1pt]$R=S\cdot C$};
  \node[card=amber, minimum width=54mm, minimum height=13mm] (G) at (3.4,-3.6) {progress\\[-1pt]$\dot S/S=R^{\lambda}S^{-\beta}$};
  \node[card=slate, minimum width=36mm, minimum height=9mm] (C) at (6.8,2.3) {compute $C$ to run the AIs};
  \draw[->, accent, line width=1pt] (S.east) -- node[above, smallnode] {smarter, cheaper AIs} (R.west);
  \draw[->, slate, line width=0.8pt] (C.south) -- (R.north);
  \draw[->, teal, line width=1pt] (R.south) -- node[right, smallnode, align=left, xshift=1mm] {more research,\\ penalty $\lambda$} (G.north east);
  \draw[->, amber!85!black, line width=1pt] (G.north west) -- node[left, smallnode, align=right, xshift=-1mm] {ideas get harder:\\ factor $S^{-\beta}$} (S.south);
\end{tikzpicture}
\caption{The feedback loop of the Jones model once AIs do the research. Better software makes the AI researchers better (and cheaper to run), which raises research input $R$; research input drives progress, with a parallelization penalty $\lambda$ on the way in and a diminishing-returns factor $S^{-\beta}$ on the way out. Everything hinges on whether the loop \emph{amplifies} ($\lambda>\beta$).}
\label{fig:L4-loop}
\end{figure}

\begin{prop}[Dynamics of the Jones model\dagmark]\label{prop:L4-loop}
Let $\lambda>0$, $S_0>0$, and let $S>0$ solve \eqref{eq:L4-jones}.
\begin{enumerate}[label=\textbf{(\alph*)}]
  \item Along every solution, $\dfrac{\dd}{\dd t}S^{\beta}=\beta R^{\lambda}$. For a given (exogenous) continuous input $R(t)>0$ and $\beta>0$ it follows that $S^{\beta}=S_0^{\beta}+\beta\int_0^tR^{\lambda}\,\dd s$, so $S$ exists for all times.
  \item Let $\beta>0$. If $R=CS$ with $C>0$ constant, then $u:=S^{\beta}$ solves $\dot u=k\,u^{\,r}$ with $k=\beta C^{\lambda}$ and $r=\lambda/\beta$. For $r>1$ it blows up at
  \[ T^*=\frac{u_0^{\,1-r}}{(r-1)\,k},\qquad u(t)=\big(u_0^{\,1-r}-(r-1)kt\big)^{-1/(r-1)}\quad(0\le t<T^*); \]
  for $r=1$ it is exponential, and for $r<1$ it exists for all times and grows polynomially, $u\asymp t^{1/(1-r)}$, so $S\asymp t^{1/(\beta(1-r))}$.
  \item Let $\beta>0$. If $R(t)=R_0e^{g_Rt}$ with $R_0>0$ and $g_R>0$, then $\dfrac1t\log S\to r\,g_R$: software quality grows exponentially at $r$ times the growth rate of research effort.
\end{enumerate}
\end{prop}
\begin{proof}
(a) $\frac{\dd}{\dd t}S^\beta=\beta S^{\beta-1}\dot S=\beta S^{\beta-1}\cdot S\,R^{\lambda}S^{-\beta}=\beta R^\lambda$. (b) Insert $R=CS=Cu^{1/\beta}$ in (a): $\dot u=\beta C^\lambda u^{\lambda/\beta}$; then separate variables. (c) By (a), $S^\beta=S_0^\beta+\beta\int_0^tR_0^\lambda e^{\lambda g_Rs}\,\dd s=S_0^\beta+\frac{\beta R_0^\lambda}{\lambda g_R}\,(e^{\lambda g_Rt}-1)$, so $\log S=\tfrac1\beta\log(\cdots)=\frac{\lambda g_R}{\beta}t+O(1)$. The sign of $g_R$ matters: for $g_R=0$ one gets $S^\beta=S_0^\beta+\beta R_0^\lambda t$ and the limit is $0=r\,g_R$, and for $g_R<0$ the integral converges, $S$ stays bounded, and the limit is $0\neq r\,g_R$.
\end{proof}
\noindent The lecture states the condition \eqref{eq:L4-r}; the closed forms and the proof are mine. Part (b) is exactly the toy ODE of Lecture~1 (Section~\ref{sec:L1-rsi}), $\dot x=kx^{p}$, with $x=S^{\beta}$ and $p=r$: the Jones model is that toy model with its exponent given an economic meaning. Part (c) says what ``returns'' means: on a balanced growth path $S\propto R^{\,r}$, so each doubling of research effort buys $r$ doublings of software quality. And $r>1$ is where ``AI helps AI research'' is strong enough to win the race against diminishing returns.

\begin{remark}[When is $r$ the right parameter?\dagmark]
The ratio $r=\lambda/\beta$ only makes sense for $\beta>0$ (ideas get harder). If $\beta\le0$ (the growth rate itself rises with $S$, which the speaker calls ``standing on the shoulders of giants''; in Jones' own vocabulary that name is used already for $\phi=1-\beta>0$, i.e.\ for $\beta<1$) then $1+\lambda-\beta>1$ for every $\lambda>0$ and the loop \eqref{eq:L4-closed} blows up whatever $\lambda$ is; the primary condition is $\lambda>\beta$, and $r>1$ is its restatement for $\beta>0$.
\end{remark}

\begin{figure}[tb]
\centering
\begin{tikzpicture}
\begin{axis}[notes, width=0.88\linewidth, height=0.5\linewidth, ymode=log, xmin=0, xmax=3, ymin=1, ymax=1e4,
   xlabel={time $t$ (units in which $k=1$ and $u_0=1$)}, ylabel={software quality $u=S^{\beta}$},
   legend style={at={(0.02,0.98)}, anchor=north west, font=\sffamily\fontsize{6.2}{7}\selectfont}]
  \addplot[accent, line width=1.6pt] table[x=t, y=r0p5] {figures/data/jones_traj.dat};   \addlegendentry{$r=0.5$ (polynomial)}
  \addplot[ink, line width=1.6pt] table[x=t, y=r1p0] {figures/data/jones_traj.dat};   \addlegendentry{$r=1$ (exponential)}
  \addplot[teal, line width=1.8pt] table[x=t, y=r1p5] {figures/data/jones_traj.dat};   \addlegendentry{$r=1.5$, $T^*=2$}
  \addplot[amber, line width=1.8pt] table[x=t, y=r2p0] {figures/data/jones_traj.dat};   \addlegendentry{$r=2$, $T^*=1$}
  \addplot[crimson, line width=1.8pt] table[x=t, y=r3p0] {figures/data/jones_traj.dat};   \addlegendentry{$r=3$, $T^*=0.5$}
  \draw[teal, dotted, line width=0.8pt] (axis cs:2,1) -- (axis cs:2,1e4);
  \draw[amber, dotted, line width=0.8pt] (axis cs:1,1) -- (axis cs:1,1e4);
  \draw[crimson, dotted, line width=0.8pt] (axis cs:0.5,1) -- (axis cs:0.5,1e4);
\end{axis}
\end{tikzpicture}
\caption{Solutions of $\dot u=k\,u^{r}$, $u=S^{\beta}$ (Proposition~\ref{prop:L4-loop}), drawn on a logarithmic scale: for $r<1$ the curve bends \emph{down} (growth slower than exponential), for $r=1$ it is a straight line, for $r>1$ it bends \emph{up} and reaches infinity at $T^*=u_0^{1-r}/((r-1)k)$ (dotted). Exact solutions, checked against numerical integration of the original equation \eqref{eq:L4-jones} with $R=CS$ (script \kbd{figures/gen\_data.py}).}
\label{fig:L4-traj}
\end{figure}

\subsection{A back-of-the-envelope estimate of $r$}\label{sec:L4-balanced}\tsm{11:55}
How does one estimate $r$? The full problem (maximum likelihood, Bayesian methods, the whole time series of $S$ and $R$) is what the speaker's colleagues do, and he does not go into it; he shows only a naive approach. Assume a \term{balanced growth path}, on which the growth rate $g_S:=\dot S/S$ of software quality is constant. Take logarithms of \eqref{eq:L4-jones} and differentiate in time:
\[
   \frac{\dot g_S}{g_S}\;=\;\lambda\,g_R-\beta\,g_S,\qquad g_R:=\frac{\dot R}{R}.
\]
Balanced growth means $\dot g_S=0$, hence
\begin{equation}\label{eq:L4-balanced}
   r\;=\;\frac{\lambda}{\beta}\;=\;\frac{g_S}{g_R}.
\end{equation}
The growth rate of the \emph{output} (software quality) divided by the growth rate of the \emph{input} (research effort) is the returns to R\&D. If one knows both growth rates in some domain, one has a back-of-the-envelope value of $r$ there. This is a one-line version of Proposition~\ref{prop:L4-loop}(c). (A caveat of mine: \eqref{eq:L4-balanced} is a calibration identity for an \emph{open} loop, in which the research input is given by data. In the closed loop $R=CS$ one has $g_R=g_S$, so a balanced path would need $r=1$; for $r>1$ the growth is hyperbolic and for $r<1$ polynomial, by Proposition~\ref{prop:L4-loop}(b).)

%% ------------------------------------------------------------------------------------
\section{Measuring software progress and estimating $r$}\label{sec:L4-measure}
\tsm{13:27}
Now one needs data, and the speaker concentrates on this part because it is what he has worked on most. The first problem: what does ``software quality'' actually mean?

\paragraph{Software quality as compute efficiency.}\tsm{13:58}
For classical algorithms one speaks of complexity, $O(n\log n)$ and so on, but for a neural network that is hard to define. The more common approach in machine learning is to ask: \emph{how much less training compute is needed to reach the same level of capability?} If reaching GPT-4-level capability now takes ten times less training compute, the algorithms and data are more efficient. One picks a reference algorithm or data family, and defines everything else relative to it, which avoids having to state the computational complexity of an arbitrary architecture.
Empirically, algorithms obey \emph{scaling relationships} between training compute (on a log scale) and capability along some metric; in the stylized picture of the lecture, capability rises linearly in $\log C$ (Lecture~2). That is why the labs buy so many GPUs. (The speaker adds that this year's spending on data centers in the US, mostly on chips, is in the hundreds of billions of dollars, and as a fraction of GDP exceeds the Manhattan Project and the Apollo program\flag.)
An efficiency gain \emph{shifts the curve to the left}: less compute for the same capability; equivalently, more capability at the same compute (Figure~\ref{fig:L4-shift}a). A student asks whether the shift must be parallel; no, ``I will get to that later'' (Section~\ref{sec:L4-scale}); the picture is stylized.

\begin{figure}[tb]
\centering
\begin{tikzpicture}
\begin{groupplot}[group style={group size=2 by 1, horizontal sep=9mm}, width=0.48\linewidth, height=0.45\linewidth,
    axis lines=left, axis line style={-{Latex[length=2mm,width=1.4mm]}, ink!70, line width=0.5pt}, ticks=none,
    xmin=0, xmax=4, ymin=0, ymax=3.5, xlabel={log training compute}, ylabel={loss},
    label style={font=\sffamily\small, text=ink}, title style={font=\sffamily\small\bfseries, text=ink, yshift=-1mm}]
\nextgroupplot[title={(a) parallel shift}]
  \addplot[slate, line width=1.6pt, domain=0:4] {3.2-0.7*x};
  \addplot[accent, line width=1.8pt, domain=0:3.4] {2.5-0.7*x};
  \draw[<->, amber!85!black, line width=1pt] (axis cs:1.43,1.5) -- (axis cs:2.43,1.5);
  \node[font=\sffamily\fontsize{6.4}{7}\selectfont, text=amber!80!black, anchor=south west, align=left, fill=white, fill opacity=0.85, text opacity=1, inner sep=1pt] at (axis cs:2.55,1.72) {same loss,\\ $G\times$ less compute};
  \node[font=\sffamily\fontsize{6.4}{7}\selectfont, text=slate, anchor=south west] at (axis cs:0.08,2.75) {old};
  \node[font=\sffamily\fontsize{6.4}{7}\selectfont, text=accent, anchor=north east] at (axis cs:1.6,0.35) {new};
\nextgroupplot[title={(b) scale-dependent shift}]
  \addplot[slate, line width=1.6pt, domain=0:4] {3.2-0.7*x};
  \addplot[accent, line width=1.8pt, domain=0:3.4] {3.45-1.0*x};
  \draw[<->, amber!85!black, line width=1pt] (axis cs:1.45,2.0) -- (axis cs:1.714,2.0);
  \draw[<->, amber!85!black, line width=1pt] (axis cs:2.75,0.7) -- (axis cs:3.571,0.7);
  \node[font=\sffamily\fontsize{6.4}{7}\selectfont, text=amber!80!black, anchor=west, fill=white, fill opacity=0.85, text opacity=1, inner sep=1pt] at (axis cs:1.85,2.0) {$G_1$ small};
  \node[font=\sffamily\fontsize{6.4}{7}\selectfont, text=amber!80!black, anchor=south, fill=white, fill opacity=0.85, text opacity=1, inner sep=1pt] at (axis cs:3.2,0.8) {$G_2$ large};
\end{groupplot}
\end{tikzpicture}
\caption{Schematic scaling curves (no data). \textbf{(a)} A constant efficiency gain translates the curve to the left by a fixed amount: the same loss needs $G$ times less compute at every scale. \textbf{(b)} Scale-dependent progress: a new method with a steeper curve saves little compute at small scale and a lot at large scale (here it is even worse at the smallest scales, where the curves cross). This is the situation of Section~\ref{sec:L4-scale}.}
\label{fig:L4-shift}
\end{figure}

\paragraph{The rate of software progress.}\tsm{17:02}
The rate of AI software progress is then how much less training compute is needed to reach a given capability level as time passes, possibly averaged over several capability levels. The difficulty: outside the labs, where everything is kept secret, there is no public data from which to get a clean time series. The rough approach: take a typical scaling law, with cross-entropy loss as a function of compute, and augment it with a \emph{software efficiency factor} $S(t)$, usually assumed to grow exponentially in time, $S(t)=e^{kt}$; the coefficient $k$ is then estimated, which gives the overall rate of software progress and the doubling time $\ln2/k$. This needs a big dataset: the training compute and the validation loss of many models.

\begin{beyond}[The model behind the number]
This is the approach of \citet{ho2024}, of which the speaker is the first author. Starting from the parametric law $L=E+AN^{-a}+BD^{-b}$ of Lecture~2, they replace the model size and the number of training tokens by \emph{effective} ones, $N_{\mathrm{eff}}=N\,e^{a'(Y-Y_0)}$ and $D_{\mathrm{eff}}=D\,e^{b'(Y-Y_0)}$ (the paper writes $\alpha,\beta,\alpha',\beta'$), with $Y$ the publication year, so that algorithmic progress is an exponential trend multiplying both terms, and effective compute follows from $C\approx6ND$. Fitted to more than 200 language-model evaluations from 2012 to 2023 (Wikitext and Penn Treebank), the compute required for a fixed level of performance halved roughly every $8$ months (95\% confidence interval about $5$ to $14$ months in the first version of the paper; the main text gives $8.4$ months with interval $4.5$ to $14.3$). Eight months is $2^{12/8}\approx2.8$ per year, which is the ``central estimate of about three times per year'' of the lecture, and the interval endpoints correspond to about $1.8\times$ and $5.3\times$ per year. These endpoints are far narrower than the ``two to forty or fifty'' per year that the speaker recalls; I could not tell whether his range refers to the spread across the several studies he plots, to a later analysis, or is a slip of memory. The same paper estimates that the transformer was worth a median compute-equivalent gain of $7.2\times$ (95\% interval $3.3\times$ to $45.7\times$) over earlier architectures (LSTMs and others) at a frontier budget of $10^{25}$ FLOP, and $6.6\times$ at $10^{22}$ FLOP, and cautions that gains depend on scale (Section~\ref{sec:L4-scale}); these numbers are not those of the slide discussed there, and I do not know which comparison that slide shows.
\end{beyond}

\paragraph{Why the error bars are huge.}\tsm{18:37}
Because the data are poor, the intervals are enormous. The original estimate, with a central value of about $3\times$ per year, has, as he reads off the upper and lower estimates of the analysis, a range from ``maybe two'' to something like forty or fifty times per year (he does not remember the numbers)\flag. He then shows a graph that puts several estimates from the literature side by side: if one believes them, a central estimate of order $10\times$ per year is plausible, i.e.\ ten times less training compute for the same capability each year. The methodologies differ (for instance, in whether they count distillation or synthetic-data generation), and in every case the uncertainty bounds are wide.

\begin{caution}[What the estimates do not cover]
\begin{itemize}
  \item They concern \textbf{pre-training}. Nearly all models in use today are \emph{reasoning models} trained with much reinforcement learning (``the models that help crack the Jacobian conjecture or the Navier--Stokes problem''), and their training compute is unknown; yet compute is needed to attribute progress to software. So the datasets do not capture this approach.
  \item Progress may be \textbf{faster inside the labs}. The model said to have solved the Navier--Stokes problem was reportedly an internal model that did so within about three days, before its training had even finished\flag; it is far more capable than what is available outside. Perhaps software progress is just steeper there (``a shift in the slope''); from outside, one cannot tell.
\end{itemize}
\end{caution}

\paragraph{Estimating $r$ itself.}\tsm{22:16}
Even so, given the rough growth rate of $S$ one can try to estimate $r$ with \eqref{eq:L4-balanced} by also getting the growth rate of the researchers, using proxies for research input: the number of papers published in a field, or the number of researchers at OpenAI over time. The result is again ``kind of ridiculous'': middle estimates are sometimes above one, but one cannot rule out being below or above. And those estimates used the old $3\times$ per year (the 2024 estimate); since $r$ is the ratio of the growth rate of software to that of the researchers, a faster $g_S$ means a proportionally larger $r$. ``So naively this suggests that something like a software intelligence explosion is plausible, although there are many issues, and I personally do not put much weight on it.''

\begin{beyond}[A correction in the arithmetic, and the economy-wide benchmark]
The speaker says that moving from $3\times$ to $10\times$ per year ``might increase the returns by something like three times''. In terms of growth \emph{rates}, $g_S=\ln(\text{annual factor})$, so the increase is $\ln10/\ln3\approx2.1$, a factor of two rather than three; the point, that the naive $r=g_S/g_R$ is proportional to $g_S$, stands. Epoch's own estimates of $r$ for software domains (\citealp{besiroglu2024}; the paper behind it is ``Estimating idea production'' by Erdil, Besiroglu and Ho) use a Bayesian update with the number of authors as the proxy for research input: for computer chess (Stockfish) they find a central value of $r\approx0.83$ with standard error $0.15$, and for computer vision, SAT solvers, linear programming and reinforcement learning median estimates above $1$ that are \emph{not} statistically significant. Their replication of the economy-wide estimate of \citet{bloom2020} gives about $0.25$ against the original $0.32$ (the ``$0.3$'' of the lecture). I take these numbers from Epoch's post.
\end{beyond}

\begin{discussion}[Does $S$ depend on the training compute?]
A student notices that in $R=S\cdot C$ the efficiency $S$ is defined through training compute, and asks how training compute is related to $r$. The speaker: the right way to think is \emph{not} through that $S\cdot C$ model, which is a base model that is ``clearly incorrect conceptually in some ways'' and works only under assumptions (Section~\ref{sec:L4-two}).
\end{discussion}

% =====================================================================================
%  Lecture 4, part B  —  section 4.4 (extensions), 4.5 (questions) and closing material
% =====================================================================================

%% ------------------------------------------------------------------------------------
\section{Beyond the basic model}\label{sec:L4-extend}
\tsm{24:22}
The speaker now builds on the basic picture and takes up the questions that have been raised. The biggest problems with the basic model, he says, are that it \emph{misses some dynamics}, such as the dependence of algorithmic efficiency on the scale of the training compute, and that it \emph{misses inputs} to the production function of R\&D. Four extensions follow.

\subsection{Experiments need compute: the CES production function}\label{sec:L4-ces}\tsm{24:54}
To discover new algorithms and make more software progress one may have to \emph{run many experiments} on physical compute: ``it is not purely a case of being smart; no matter how smart you are, maybe you just need to test things out''. If so, the research input has to include a stock of \emph{experiment compute} besides the cognitive labor of the AIs. A paper by Parker Whitfill and Cheryl Wu\,\citep{whitfill2025} tries to quantify these bottlenecks by defining $R$ as a \term{CES} (constant elasticity of substitution) function of compute and cognitive labor. At the board the speaker explains what that is (``is there a whiteboard?''); he notes that mathematicians have a similar object, the $\ell^p$ norm, but that here the exponent ranges over a different set of values.

\needspace{9\baselineskip}
\begin{defn}[CES production function]\label{def:L4-ces}
Let $K,L>0$ be two inputs (economists: capital and labor; here: experiment compute $K$ and cognitive labor $L$). For a weight $\alpha\in(0,1)$ and a substitution parameter $\rho\le1$,
\[
   R(K,L)\;=\;\big(\alpha\,K^{\rho}+(1-\alpha)\,L^{\rho}\big)^{1/\rho}\quad(\rho\ne0),\qquad R(K,L)=K^{\alpha}L^{1-\alpha}\quad(\rho=0).
\]
\end{defn}
\noindent It is a weighted \emph{power mean}, one family that contains the arithmetic, geometric and harmonic means (my remark: the range $\rho\le1$ is the one economists use; it keeps the level curves convex). The speaker walks through the limits (Table~\ref{tab:L4-ces}): for $\rho=1$ the two inputs are \emph{perfect substitutes}, ``if you increase this, then the whole thing also increases''; as $\rho\to-\infty$ the function tends to $\min(K,L)$, the inputs are \emph{complements}: if one of them is $1000$ and the other $1$, no amount of increasing the big one raises the output much, ``you are bounded by the other thing''. The special case $\rho=0$, by l'H\^opital's rule, is Cobb--Douglas, $K^{\alpha}L^{1-\alpha}$ (``left as an exercise for the reader'').
The economist's examples: with $\rho=-\infty$, output is the number of \emph{pairs of shoes}, with $K$ left shoes and $L$ right shoes (a thousand left shoes and one right shoe still make one pair); for substitutes, output is \emph{fruit} and it does not matter whether one has apples or pears.

\begin{table}[tb]
\centering\small
\renewcommand{\arraystretch}{1.28}
\begin{tabular}{@{}>{\raggedright\arraybackslash}p{17mm}>{\raggedright\arraybackslash}p{34mm}>{\raggedright\arraybackslash}p{44mm}>{\raggedright\arraybackslash}p{16mm}@{}}
\toprule
\textbf{$\rho$} & \textbf{$R(K,L)$} & \textbf{Name; bottleneck?} & \textbf{$\sigma=\frac1{1-\rho}$}\\
\midrule
$1$ & $\alpha K+(1-\alpha)L$ & perfect substitutes; none & $\infty$\\
$0<\rho<1$ & power mean & substitutes; none ($R\to\infty$ as $L\to\infty$) & $>1$\\
$\rho\to0$ & $K^{\alpha}L^{1-\alpha}$ & Cobb--Douglas; $R\to\infty$, but slowly & $1$\\
$\rho<0$ & power mean & complements; \textbf{ceiling} $R<\alpha^{1/\rho}K$ & $<1$\\
$\rho\to-\infty$ & $\min(K,L)$ & perfect complements & $0$\\
\bottomrule
\end{tabular}
\caption{The CES family. The last column is the \emph{elasticity of substitution} that gives the family its name: $\sigma=\dd\ln(L/K)/\dd\ln\mathrm{MRS}$ along a level curve, where $\mathrm{MRS}=F_K/F_L$ is the ratio of the marginal products (the slope $\abs{\dd L/\dd K}$ of the level curve). Substitutes ($\sigma>1$) and complements ($\sigma<1$) are separated by Cobb--Douglas \citep{arrow1961}. The $\sigma$ column and the ceilings are my additions\dagmark{}.}
\label{tab:L4-ces}
\end{table}

\paragraph{What it means for an explosion.}\tsm{29:03}
In the application $K$ is the experiment compute and $L$ the cognitive labor. If $\rho<0$ the two are complements, ``like the left shoes and the right shoes'', and one becomes a \emph{bottleneck} for the other: one may increase the intelligence, or the number, of the AI systems a lot, but one is then bottlenecked on experiments. If $\rho>0$ they are substitutes, there is no bottleneck, and ``you can increase cognitive labor to infinity and you are golden: you get the hyperbolic growth path'' (in the model of Section~\ref{sec:L4-jones} this needs $r>1$ as well; the lecture leaves that implicit). A student remarks (the captions are garbled\flag) that for $\rho<0$ the output simply converges to a finite limit as $L\to\infty$; ``exactly''.

\begin{figure}[tb]
\centering
\begin{tikzpicture}
\begin{groupplot}[group style={group size=2 by 1, horizontal sep=11mm}, width=0.5\linewidth, height=0.5\linewidth,
    axis lines=left, axis line style={-{Latex[length=2mm,width=1.4mm]}, ink!70, line width=0.5pt},
    tick style={ink!70, line width=0.4pt}, every tick label/.append style={font=\sffamily\scriptsize, text=ink!85},
    label style={font=\sffamily\small, text=ink}, grid=major, major grid style={rule!70, line width=0.3pt},
    every axis plot/.append style={line width=1.4pt}, unbounded coords=jump, table/col sep=space,
    title style={font=\sffamily\small\bfseries, text=ink, yshift=-1mm}]
\nextgroupplot[title={(a) output with $K=1$}, xmode=log, xmin=0.1, xmax=1000, ymin=0, ymax=6.4, xlabel={cognitive labor $L$}, ylabel={output $R(1,L)$},
   legend to name=cesleg, legend columns=-1, legend style={font=\sffamily\fontsize{6.4}{7}\selectfont, draw=none, /tikz/every even column/.append style={column sep=3.5mm}}]
  \addplot[crimson] table[x=L, y=rho1] {figures/data/ces_labor.dat};      \addlegendentry{$\rho=1$}
  \addplot[amber] table[x=L, y=rho0p5] {figures/data/ces_labor.dat};      \addlegendentry{$\rho=0.5$}
  \addplot[moss] table[x=L, y=rho0] {figures/data/ces_labor.dat};         \addlegendentry{$\rho=0$}
  \addplot[teal] table[x=L, y=rhom0p4] {figures/data/ces_labor.dat};      \addlegendentry{$\rho=-0.4$}
  \addplot[accent] table[x=L, y=rhom1] {figures/data/ces_labor.dat};      \addlegendentry{$\rho=-1$}
  \addplot[plum] table[x=L, y=rhom4] {figures/data/ces_labor.dat};        \addlegendentry{$\rho=-4$}
  \draw[teal, dashed, line width=0.7pt] (axis cs:0.1,5.657) -- (axis cs:1000,5.657);
  \draw[accent, dashed, line width=0.7pt] (axis cs:0.1,2) -- (axis cs:1000,2);
  \draw[plum, dashed, line width=0.7pt] (axis cs:0.1,1.189) -- (axis cs:1000,1.189);
  \node[font=\sffamily\fontsize{5.8}{6.4}\selectfont, text=teal, anchor=south east, fill=white, fill opacity=0.85, text opacity=1, inner sep=1pt] at (axis cs:1000,5.72) {ceiling $2^{2.5}\approx5.7$ for $\rho=-0.4$};
\nextgroupplot[title={(b) level curves $R=1$}, xmin=0, xmax=2.4, ymin=0, ymax=2.4, xlabel={experiment compute $K$}, ylabel={cognitive labor $L$}]
  \addplot[crimson] table[x=K, y=rho1] {figures/data/ces_isoquant.dat};
  \addplot[moss] table[x=K, y=rho0] {figures/data/ces_isoquant.dat};
  \addplot[accent] table[x=K, y=rhom1] {figures/data/ces_isoquant.dat};
  \addplot[plum] table[x=K, y=rhom4] {figures/data/ces_isoquant.dat};
\end{groupplot}
\end{tikzpicture}\par\smallskip\ref{cesleg}
\caption{The CES function with equal weights $\alpha=\tfrac12$ (computed from the formula; script \kbd{figures/gen\_data.py}). \textbf{(a)} Output as the cognitive labor grows with experiment compute fixed at $K=1$. Substitutes ($\rho>0$) have no ceiling; for $\rho<0$ the output saturates at $\alpha^{1/\rho}K$ (dashed), and for $\rho$ close to $0$ the ceiling is approached very slowly. \textbf{(b)} Level curves $R=1$ (colors as in the legend, for $\rho=1,0,-1,-4$): for complements they become an L-shape with asymptotes at $K=\alpha^{1/|\rho|}$ and $L=(1-\alpha)^{1/|\rho|}$; no amount of the plentiful input compensates for a shortage of the other.}
\label{fig:L4-ces}
\end{figure}

\paragraph{What the data say.}\tsm{30:05}
In the \emph{baseline} model of Whitfill and Wu, with total experiment compute and cognitive labor as the inputs, the bottleneck is quite weak: they find $\rho>0$. If instead of the total stock of experiment compute one uses the number of \emph{frontier experiments}, i.e.\ experiments at or near the scale of the largest training runs, then the bottleneck becomes strong and $\rho$ is very negative.
A student asks what ``frontier'' means here. Two pictures: one can run huge numbers of tiny experiments, at the scale of GPT-2, and hope that the innovations transfer to the largest runs. But people at the frontier labs point out that one does not know whether an algorithm \emph{scales}: it may work at the GPT-2 scale and break at the GPT-6 scale, giving no efficiency boost. So one may need to run experiments across scales and close to the scale where the algorithm is to be used. In the model, the total R\&D compute stock is divided by the size of the training run, with a factor that allows generalization two or three orders of magnitude away from the frontier; with that adjustment the compute bottleneck becomes much stronger. Noam Brown\flag{} of OpenAI is cited as saying that many algorithmic innovations introduced in academia do not work at frontier scales.

\begin{discussion}[Is there evidence for the Jones model itself?]
A student finds it odd that the exponent $\lambda$ does not depend on $S$ (the same ``difficulty'' at all levels of knowledge) and asks for the evidence for the formula. The speaker: evidence for it is probably not available in the AI world. It comes from \emph{endogenous growth theory}, where Jones introduced the exponents to repair an unrealistic prediction of earlier models, in which raising research effort gives a \emph{permanent increase in the growth rate} (twice the researchers, twice the growth rate, forever); he is not sure where exactly the exponent $\beta$ comes from, and will look it up.
\end{discussion}

\begin{beyond}[Why Jones needed the exponents, and what $r$ was there]
In Romer's model\,\citep{romer1990} new ideas are produced by $\dot A=\delta L_AA$, so the growth rate $\dot A/A=\delta L_A$ is proportional to the \emph{number} of researchers: a ``scale effect''. \citet{jones1995} replaced it by $\dot A=\delta L_A^{\lambda}A^{\phi}$ with $\phi<1$, which is \eqref{eq:L4-jones} with $\beta=1-\phi$. If research labor grows at a constant rate $n$, the long-run growth rate of $A$ is $\lambda n/(1-\phi)=r\,n$, exactly the relation \eqref{eq:L4-balanced}: the growth rate of technology is tied to the \emph{growth rate} of the research population, not to its level (``semi-endogenous'' growth). So $r$ is the multiplier from research growth to idea growth, and Section~\ref{sec:L4-balanced} is Jones' formula read backwards.
\end{beyond}

\begin{beyond}[What a bottleneck does to the explosion condition]\label{box:L4-bottleneck}
Combine the two models, for $\beta>0$: $R=\mathrm{CES}(K,L)$ with experiment compute $K$ and cognitive labor $L=CS$ as in \eqref{eq:L4-RSC}, inserted in \eqref{eq:L4-jones}. Proposition~\ref{prop:L4-loop}(a) still holds, $\frac{\dd}{\dd t}S^\beta=\beta R^\lambda$, and the behavior of $R$ for large $L$ decides everything. Write $R\asymp L^{\theta}$ for large $L$, with $K$ fixed: $\theta=1$ if $0<\rho\le1$, $\theta=1-\alpha$ if $\rho=0$, and $\theta=0$ if $\rho<0$.
\begin{enumerate}[label=(\roman*)]
  \item \textbf{$K$ constant.} The loop blows up in finite time iff $\theta\,r>1$. That is: for substitutes iff $r>1$ (the basic condition); for Cobb--Douglas iff $(1-\alpha)\,r>1$; for complements \emph{never}, and then $S^\beta$ grows at most linearly, $S\lesssim t^{1/\beta}$.
  \item \textbf{$\rho<0$ and $K(t)=K_0e^{g_Kt}$, $g_K\ge0$.} Then $\limsup_{t\to\infty}\frac1t\log S\le r\,g_K$, with equality if $r>1$; if $r<1$ the limit is $0$ (AI labor, not compute, is then the limiting factor: polynomial growth). So growth is at most exponential, at a rate set by the growth of \emph{compute}, not of AI labor.
\end{enumerate}
\emph{Proof.} For $\rho<0$ the map $x\mapsto x^{1/\rho}$ is decreasing, so $R\le(\alpha K^{\rho})^{1/\rho}=\alpha^{1/\rho}K$, and (a) gives $S^{\beta}\le S_0^\beta+\beta\alpha^{\lambda/\rho}K^\lambda t$ for constant $K$, and $S^\beta\le S_0^\beta+\beta\alpha^{\lambda/\rho}K_0^\lambda(e^{\lambda g_Kt}-1)/(\lambda g_K)$ in case (ii), i.e.\ $\log S\le r g_Kt+O(1)$. For $\rho=0$, $R=K^\alpha(CS)^{1-\alpha}$, so $u=S^\beta$ solves $\dot u=k'u^{(1-\alpha)r}$: the equation of Proposition~\ref{prop:L4-loop}(b) with $r$ replaced by $(1-\alpha)r$. For $0<\rho\le1$ one has $(1-\alpha)^{1/\rho}CS\le R\le K+CS$ (a weighted power mean never exceeds the largest input), which gives $\dot u\ge c\,u^{r}$ (blow-up if $r>1$) and, from the upper bound, at most exponential growth if $r\le1$.

For the remaining claims of (ii), note that $R\le(1-\alpha)^{1/\rho}CS$ as well (same argument), so $\dot u\le c_1u^r$ with $c_1=\beta\big((1-\alpha)^{1/\rho}C\big)^{\lambda}$: for $r<1$ the quantity $u$ grows at most polynomially and $\log S=o(t)$. For $r>1$ the upper bound is attained. A power mean is at least the smaller input, $R\ge\min(K,CS)$. Let $u_*=(K/C)^\beta$, so that $CS\ge K$ iff $u\ge u_*$. If $u<u_*$ for all large $t$, then $\dot u\ge\beta C^\lambda u^r$ and $u$ would blow up in finite time, which is absurd. At a time with $u=u_*$ one has $\dot u\ge\beta K^\lambda\propto e^{\lambda g_Kt}$, while $\dot u_*=\beta g_Ku_*\propto e^{\beta g_Kt}$ and $\lambda>\beta$; so after some time $t_c$ the difference $u-u_*$ cannot cross zero downward, and since $u\ge u_*$ at arbitrarily large times, $u\ge u_*$ from some time on. Then $\dot u\ge\beta K^\lambda$, i.e.\ $\log S\ge r\,g_Kt-O(1)$. \hfill\qedsymbol\par
The box quantifies ``bottlenecked on experiments'': once $L\gg K$, complementarity does not merely slow the explosion, it \emph{switches the feedback off} (the exponent $\theta=0$), and what remains is a growth rate bounded by the growth of hardware (the speaker cites chips below); while $L\ll K$ the loop still runs as in Proposition~\ref{prop:L4-loop}. The dichotomy is asymptotic ($L\to\infty$) and not uniform in $\rho$: with equal weights the ceiling is $2^{1/|\rho|}K$, i.e.\ $5.7K$ at $\rho=-0.4$ but $1024K$ at $\rho=-0.1$ and $10^6K$ at $\rho=-0.05$, and as $\rho\downarrow0$ the blow-up time in the substitutes case diverges (for $\lambda=1.5$, $\beta=1$, $\alpha=\tfrac12$, $K=C=S_0=1$: $T^*\approx6,\,13,\,30,\,104,\,2150$ for $\rho=0.5,\,0.2,\,0.1,\,0.05,\,0.02$). So the sign of $\rho$ decides whether an explosion is possible in principle, and its size decides on what horizon. I checked all of this numerically (script \kbd{figures/check\_L4\_bottleneck.py}: the blow-up classification for $\rho\in\{1,0.5,0,-0.5,-2\}$ and $r\in\{0.5,1,1.5,1.9,2.1\}$, the growth rates in (ii), and the numbers above).
\emph{How large is the ceiling?} For equal weights and $\rho=-0.4$ the ceiling $\alpha^{1/\rho}K$ is $2^{2.5}\approx5.7$ times the output $R(K,K)=K$ one has when labor equals compute. This is consistent with the ``maximum speed-up of about $6\times$'' that \citet{davidson2025} derives, in his essay on this very objection, from an estimate $\rho\approx-0.4$ for US manufacturing (an elasticity of substitution of about $0.7$); he argues that for AI R\&D $\rho$ is probably higher (between $-0.2$ and $0$) and gives a 10--40\% chance of an explosion despite compute bottlenecks. (I compared these numbers with the passages of the essay, read through a page-summarizing tool; check the essay itself before citing.)
\end{beyond}

\subsection{Intelligence versus the number of AIs}\label{sec:L4-two}\tsm{34:45}
A second issue is that the model does not cleanly separate \emph{improvements in intelligence} from \emph{increases in the number of AIs}. In $R=S\cdot C$ the speaker has been taking $S$ to be training-compute efficiency, while $C$ is more like inference compute, so multiplying them is ``sort of a category error''; it is justified only if training efficiency and inference efficiency are (linearly) related, in which case the setup works.
One way to separate the two is the AI Futures model\,\citep{aifutures2025dec}, ``probably the best model of the software intelligence explosion to date'', because it accounts for all this in much more depth; a student asks whether it is the AI~2027 model, and the answer is that it is the update of it, released in December of the previous year. In it, \emph{training compute} raises the capabilities of the models through scaling laws, which the authors relate to an abstract quantity called \emph{research taste} (how well one can use experiments); \emph{inference compute} is a separate input that goes into the amount of research effort; and experiment throughput is multiplied by research taste. ``It is pretty complicated.'' There is always a trade-off between capturing dynamics and how many parameters one can actually estimate, versus how much one is making up.

\begin{beyond}[If the two efficiencies differ, and what the AI Futures write-up says]
Suppose the AIs' efficiency at inference scales as a power $S^{\kappa}$ of the training-efficiency index: $R=S^{\kappa}C$. Then \eqref{eq:L4-closed} becomes $\dot S=C^\lambda S^{1+\kappa\lambda-\beta}$, and the explosion condition is $\kappa r>1$: the linear case $\kappa=1$ is the one in which $r$ is the right number (Exercise~7 below). As I read the write-up of the December 2025 update\,\citep{aifutures2025dec}, it is organized in three stages (automated coding, then a superhuman AI researcher, then the explosion), models experiment compute as a bottleneck on how much coding automation speeds up R\&D, and says that whether a ``taste-only'' singularity occurs depends on how fast ideas get harder to find compared with how fast research taste improves; its median parameters push the dates of full coding automation several years later than the AI~2027 model did. I have not examined the model itself.
\end{beyond}

\subsection{Scale-dependent progress}\label{sec:L4-scale}\tsm{36:50}
Another assumption is that $S$ is a \emph{single multiplier}, the same at every scale of training compute. In Figure~\ref{fig:L4-shift}(a) the curve was shifted in parallel, so that the horizontal compute-efficiency gains are always of the same size. It is possible instead that the \emph{slope} changes, and then models trained with more compute get larger efficiency gains than those trained with less (more at GPT-4 scale than at GPT-2 scale); or the slope could decrease; ``the lines may not even be straight''. The speaker simplifies in order to keep the conceptual clarity.
One case in which this has happened is the switch from recurrent networks (LSTMs) to the transformer: with the capability axis flipped to loss, the two scaling curves can be compared, and the compute saved is smaller at small scales than at large ones. The speaker quotes about $6.3\times$ near $3\cdot10^{15}$ FLOP and about $26\times$ near $10^{17}$ FLOP\flag{} (the figures are read off his slide). More generally, perhaps one cannot capture capabilities, or software quality, on a single axis at all; this is a problem regardless of scale dependence, and a more general model should allow it.

\begin{beyond}[A one-line model of scale-dependent gains]
Let the reducible losses of two methods be power laws, $L_i-E=A_iC^{-a_i}$ ($i=1$ the old method, $i=2$ the new one). To reach the reducible-loss level $\ell=L-E$ one needs $C_i=(A_i/\ell)^{1/a_i}$, so the compute-equivalent gain is $G=C_1/C_2$, and expressed through the compute $C_2$ of the new method
\[
   G\;=\;\Big(\frac{A_1}{A_2}\Big)^{1/a_1}\,C_2^{\;a_2/a_1-1}.
\]
The gain is independent of scale exactly when $a_1=a_2$ (a parallel shift of the curves on a log--log scale); if the new method has the steeper law, $a_2>a_1$, it grows like a power of the scale. The two numbers above would correspond to an exponent $\ln(26/6.3)\big/\ln\!\big(10^{17}/(3\cdot10^{15})\big)\approx0.40$ in $G\propto C_2^{\,a_2/a_1-1}$, i.e.\ $a_2/a_1\approx1.4$, if both curves were pure power laws with a common $E$ in this range and if the quoted FLOP figures are the \emph{new} method's compute $C_2$ (if they are the old method's compute $C_1=GC_2$, the same numbers give $a_2/a_1\approx1.7$); two points cannot establish any of this. \citet{ho2024} flag the same issue: older architectures such as the LSTM can be more efficient than the transformer at small scales (the curves of Figure~\ref{fig:L4-shift}b cross), and extrapolating a constant doubling time to much larger efficiency gains is unwarranted.
\end{beyond}

\subsection{Data from inside the labs}\label{sec:L4-data}\tsm{39:30}
A line of work that follows up on the basic model is to \emph{gather empirical data}, both to inform the parameters of the model and to obtain \emph{indicators}: perhaps we do not understand the model well, but we may be able to see when an explosion is under way, if it happens. People at the frontier labs, Anthropic among them, have begun to share more data about their internal R\&D process and about capability improvements inside the labs. The speaker shows a graph of the fraction of tasks done by Claude\flag{} (the captions say ``flawed'') rather than by humans, at different \emph{levels of automation} on a somewhat arbitrary scale from $0$ to $5$. Level $4$ is rising fast; a ``bold, heroic extrapolation'' has it cover almost all tasks sometime in 2027 (a student asks when; ``I was just eyeballing it''), and there is no level $5$ in the plot, except ``if you squint really hard'' at the bottom line.

This is the line of work the speaker is most excited about: without knowing many details of how AI R\&D works, it is hard to pin down the right production function, and sharing data on when an explosion might happen is extremely useful. It is also important \emph{for society} to know whether in three months ``we are going to get crazy superintelligence''. And it is important for those outside the labs to specify which metrics matter most. Epoch is actively working with labs (the speaker names Anthropic and OpenAI) on this: defining a set of \term{leading indicators} (is the Epoch Capabilities Index increasing by this much?), statistical tests that say that an acceleration of such a size triggers a particular warning, other indicators, and data to pin down the most important parameter, the complementarity between experiment compute and cognitive labor. The labs also need \emph{incentives} to share, and third-party assessments to keep them honest: METR, ``famous for looking into the Hugging Face hack'' (Lecture~1), and one other organization the captions call ``Essenture''\flag{}. There are active efforts, including at METR, at \emph{embedded investigations inside labs}, to learn about uplift and capability improvements over time; ``if we can get this out, I think that would help a lot''.

%% ------------------------------------------------------------------------------------
\section{Take-aways and questions}\label{sec:L4-questions}
\tsm{42:38}
The speaker's summary slide: (1)~the basic model of recursive self-improvement is a growth-theory model that captures ideas getting harder to find and parallelization penalties; (2)~its most important parameter is $r$, the returns to AI software R\&D, and the hyperbolic-growth condition is $r>1$; (3)~empirical estimates are terrible but do not rule out a software intelligence explosion; (4)~the basic model can be extended in several directions: experiment-compute bottlenecks, the difference between more intelligence and more workers, scale-dependent progress; and (5)~the thing he is most excited about is \emph{sharing data} on recursive self-improvement with the world.

\begin{discussion}[Credences]\label{disc:L4-credence}
\emph{Is it likely that ideas keep getting harder to find, in AI in particular, and is there reason to believe that more intelligence also gives more workers?} (The question is only partly audible.) The speaker: ``people in this crowd are supposed to be Bayesians'', so he gives a number: his personal credence in a software intelligence explosion is \textbf{about 20--25\%}. He thinks the compute bottlenecks will be pretty important and binding, but the probability is high enough that, if it happens, it could be really crazy. His best guess of $r$, which depends on the exact setup, is \emph{a little above one}. It is hard to know without a good time series of algorithmic progress, which is why Epoch asks the labs for it, or at least to estimate $r$ themselves and report it.


Asked again (and after correcting himself, having first flipped it): ideas \emph{do} get harder to find, but one can outpace that with many more or much smarter AI researchers; the 25\% is the chance that they do \emph{not} get harder to find fast enough, i.e.\ the chance of a hyperbolic path.


\emph{Is this the most widely held view?} In the San Francisco Bay Area crowd, he says, views range from ``it will not happen'' to ``it will totally happen, 100\%''; there is more variation there than in the rest of the world, so it is hard to say what the median is. His guess is that people who are really worried about AI safety would put a higher credence than he does.
\end{discussion}

\begin{discussion}[Bottlenecks in the physical world]
\emph{What are the bottlenecks for compute?} They are easier to pin down, the speaker says, because they involve things in the physical world. One of the biggest is how fast new chips can be built: look at TSMC and how fast it can expand its fabrication plants. At present it can probably sustain the current growth of compute, which he puts, from memory, at about three times a year in the number of chips\flag, but it is hard to make that go \emph{hyperbolic}; if anything it is getting harder to sustain as capital expenditure already reaches hundreds of billions of dollars; unless one improves the number of FLOP per second per chip (possible; he does not know how much AI has sped that up and guesses not much so far).


\emph{Has much work been done to estimate the substitution parameter?} Not much: the Whitfill--Wu paper is the main one, it uses basically all the public data, and everything else rests on anecdotes (``did Noam Brown think they were bottlenecked?''), which is not good for estimating $\rho$. The bottleneck is data.
\end{discussion}

\begin{discussion}[Is the capability axis unbounded? On measuring capabilities]
A student presses a deeper point. ``Less compute for the same performance'' is easy to quantify at \emph{already achieved} performance, but is there not a moment where methods stop working properly at scale? Is one assuming that capabilities form an axis that continues essentially unboundedly? (For humans one can see IQ 80 versus 140, but what is IQ 400?)


The speaker: measuring capabilities well is \emph{not a solved problem}. In the AI Futures model intelligence is captured as research taste, but nobody knows whether that is unbounded, since we have no variation in the data to see what happens if one gets there; ``it is in fact just a guess'', and a skeptic should lower the credences accordingly. One can study cases of wildly superhuman AI, chess and Go, to see how far they went before plateauing in Elo, or try to estimate the maximum software efficiency of neural networks; at the moment there is no good answer. 

It is even worse, he adds, because capabilities are not well measured \emph{today}: Anthropic's frontier safety framework (``responsible scaling policy'') asks, for instance, whether a model is at a dangerous level of biosecurity capability, requiring different safeguards or no release. At some point the models became so good that all the biosecurity benchmarks \emph{saturated}, and a saturated benchmark does not tell one whether the model could really help create a bioweapon; if it \emph{cannot} do the benchmark one is pretty confident that it cannot. It \emph{rules out} rather than \emph{rules in}.


A student suggests chess-like Elo ratings as a measurable scale. The speaker: Epoch tries several things, including Elo-like ones, but one always needs some game or common task to test on, and it need not be valid for the real world (``maybe it is just this math thing and does not apply to doing all of my job''); the Epoch Capabilities Index does this by linking overlapping benchmarks. 

And whenever a new benchmark is released, as soon as the AIs reach about $3\%$ ``you can be pretty confident that within the next two years it will be at 90-something percent, probably more like one''. For FrontierMath, Epoch interviewed Fields medalists including Terence Tao and Timothy Gowers and asked how long hard problems would last; Tao, he says, answered at least five years unless there was something like a Manhattan Project. ``Now they have blasted through the entire benchmark'', and through the harder version and some released open problems; given the spending, ``we have sort of been doing a Manhattan Project''. It is really hard to keep finding a harder benchmark that persists.
\end{discussion}

\begin{beyond}[Why a saturated benchmark can only rule things out]
Let $X$ be the event that the model has a dangerous capability and $B$ that it passes a test designed to be passed by any model with that capability: $P(B\mid X)=1$, while $P(B\mid\neg X)=q$ measures how easy the test is for models \emph{without} it. Passing multiplies the odds of $X$ by $1/q$, which tends to $1$ as the benchmark saturates ($q\to1$): no information. Failing multiplies the odds by $0$ (for $q<1$): it excludes $X$ completely. A benchmark that tests a \emph{necessary} condition therefore provides one-sided evidence whose strength is limited by how many models without the capability also pass; this is the formal content of ``rules out rather than rules in''. (If $P(B\mid X)<1$ the failure is only evidence of absence: it multiplies the odds by $(1-P(B\mid X))/(1-q)$.)
\end{beyond}

%% ------------------------------------------------------------------------------------
\tocsection{Take-aways and exercises}
\begin{openq}[Questions this lecture leaves open]
\begin{itemize}
  \item What is the right \emph{production function} of AI R\&D? Is software quality a single number at all, and does the cost of the next idea depend on the level reached, as in $S^{-\beta}$, or on the scale of the training runs?
  \item Can $r$ and $\rho$ be estimated from data that the labs could share, with error bars narrow enough to decide whether $\theta r>1$? Which \emph{leading indicators} would warn of an acceleration early enough to act?
  \item How should one model capabilities that are not bounded by any benchmark, and what happens to the feedback loop when the capability axis itself is poorly measured?
\end{itemize}
\end{openq}

\begin{takeaways}
\begin{itemize}
  \item The \textbf{software intelligence explosion} question is whether AIs doing AI research can raise capability using only software. The simplest model is Jones' growth model, $\dot S/S=R^{\lambda}S^{-\beta}$: $\lambda$ is a parallelization penalty and $\beta>0$ means ideas get harder to find.
  \item With AIs doing the research, $R=SC$ and \emph{hyperbolic growth} (a finite-time blow-up) occurs iff $\lambda>\beta$, i.e.\ $r=\lambda/\beta>1$ (for $\beta>0$). With $u=S^\beta$ it is Lecture~1's toy ODE $\dot u=ku^r$; on a balanced path $r=g_S/g_R$.
  \item \textbf{Software quality} is measured as \emph{compute efficiency}: how much less training compute reaches the same capability. Estimates of the rate ($\approx3\times$ per year in the 2024 estimate, with huge error bars) concern pre-training only; those of $r$ straddle $1$.
  \item \textbf{Compute bottlenecks.} Modeling research input as CES$(K,L)$: for complements ($\rho<0$) the output has a ceiling $\alpha^{1/\rho}K$ and once labor is plentiful the feedback switches off (for constant $K$, $S^\beta$ grows at most linearly; for $K=K_0e^{g_Kt}$, $\log S\le r\,g_Kt+O(1)$, with equality in the limit when $r>1$); Cobb--Douglas needs $(1-\alpha)r>1$; substitutes need only $r>1$ (constant $K$). Estimates of $\rho$ differ by specification: substitutes for total compute, strong complements for \emph{frontier-scale} experiments.
  \item Further refinements: separate the number of AIs from their intelligence (training versus inference efficiency); allow \emph{scale-dependent} efficiency gains; share data from inside the labs, with third-party checks.
  \item The speaker's credence for an explosion is about 20--25\% (best guess of $r$ a little above $1$); benchmark saturation means measurements of capability \emph{rule out} more than they \emph{rule in}.
\end{itemize}
\end{takeaways}

\begin{exercise}[Problems]
\begin{enumerate}
  \item \emph{The loop.} Prove parts (a) and (b) of Proposition~\ref{prop:L4-loop}. For $\lambda=0.8$, $\beta=0.5$, $C=2$, $S_0=1$ find $r$ and the blow-up time $T^*$.
  \item \emph{The polynomial regime.} For $r<1$ show that $S(t)\sim\big((1-r)kt\big)^{1/(\beta(1-r))}$, and that the solution is \emph{not} exponential.
  \item \emph{Balanced growth.} Derive \eqref{eq:L4-balanced}. If software efficiency doubles every $8$ months and the (effective) research input grows $40\%$ per year, what is $r$? What if the doubling time is $16$ months? (Hypothetical numbers.) Why does the lecture's ``three times higher'' become a factor of about two when $3\times$ per year is replaced by $10\times$ per year?
  \item \emph{CES.} Show that $R\to K^{\alpha}L^{1-\alpha}$ as $\rho\to0$ and $R\to\min(K,L)$ as $\rho\to-\infty$. Show that the elasticity of substitution is $1/(1-\rho)$.
  \item \emph{Ceilings.} For $\rho<0$ prove $R(K,L)<\alpha^{1/\rho}K$ for all $L$, and that the bound is sharp as $L\to\infty$. Evaluate it for $\alpha=\frac12$ and $\rho=-0.4,\,-1,\,-4$ (relative to $R(K,K)=K$).
  \item \emph{Bottlenecks and blow-up.} Prove the cases $\rho=0$ and $\rho<0$ of the box \emph{What a bottleneck does to the explosion condition} (Section~\ref{sec:L4-ces}). For $\rho<0$ and $K(t)=K_0e^{g_Kt}$, find the long-run growth rate of $S$ (per unit of time) when $\lambda=1.5$, $\beta=1$, $g_K=0.1$; what changes if $\lambda=0.5$?
  \item \emph{Two efficiencies.} If $R=S^\kappa C$ (inference efficiency a power of training efficiency), show that (for $\beta>0$) the loop blows up iff $\kappa r>1$, equivalently $\kappa\lambda>\beta$. What does $\kappa<1$ mean economically?
  \item \emph{Scale-dependent gains.} For $L_i-E=A_iC^{-a_i}$ derive $G=(A_1/A_2)^{1/a_1}C_2^{a_2/a_1-1}$. When does the gain not depend on the scale? What exponent do the speaker's two numbers give?
  \item \emph{Rules out, rules in.} In the model of the last Beyond box, compute the posterior probability of $X$ after passing and after failing a test with $q=0.9$, from a prior of $10\%$.
\end{enumerate}
\end{exercise}

